% Week 7 Session B — Capstone Core Development
% Overleaf: upload Week7_SessionB_Report.tex, week7_session_b_Streamlit_page.png, week7_session_b_files/
%
\documentclass[11pt,a4paper]{article}

\usepackage[margin=2.5cm]{geometry}
\usepackage{graphicx}
\newcommand{\appfig}[1]{%
  \includegraphics[width=\linewidth,height=0.62\textheight,keepaspectratio]{#1}%
}
\usepackage{booktabs}
\usepackage{xurl}
\usepackage{hyperref}
\usepackage{parskip}
\usepackage{enumitem}
\usepackage{xcolor}
\usepackage{fvextra}
\usepackage[most]{tcolorbox}

\newcommand{\studentname}{Mahmudul Hasan}
\newcommand{\studentid}{4125999049}
\newcommand{\reportdate}{May 22, 2026}
\newcommand{\capstonerepo}{https://github.com/mahmud456alhasan-debug/smart-water-capstone}

\makeatletter
\g@addto@macro\UrlBreaks{\do\ }
\makeatother

\newcommand{\LabCodeRoot}{week7_session_b_files}

\newcommand{\filebox}[2]{%
\begin{tcolorbox}[
  enhanced,
  colback=gray!6,
  colframe=black!55,
  title={#1},
  fonttitle=\bfseries\ttfamily\footnotesize,
  arc=1.5mm,
  boxrule=0.7pt,
  breakable,
  width=\linewidth,
  before skip=0.8em,
  after skip=0.8em,
]
\VerbatimInput[fontsize=\footnotesize,breaklines=true,breakanywhere=true]{\LabCodeRoot/#2}
\end{tcolorbox}%
}

\title{Lab Report: Week 7 Session B\\Capstone Core Development}
\author{\studentname \quad (Student ID: \studentid)}
\date{\reportdate}

\begin{document}
\maketitle

\begin{abstract}
This report documents Week~7 Session~B: implementation of the \textbf{Smart Water Integrated Dashboard} capstone. Core features include CSV rainfall loading with threshold alerts, SCS-CN runoff with \texttt{Q <= P} validation, seven-day reservoir dispatch (total revenue \textbf{\$708,849.26}, validation PASS), and DEM-based flood inundation maps. \texttt{pytest -q} reports \textbf{2 passed}; Streamlit runs at \texttt{localhost:8501}. Code pushed to \url{\capstonerepo} (commit \texttt{408f211}). Evidence: Figure~\ref{fig:streamlit} and appendix \texttt{prompt\_log.md}. OpenCode was not used.
\end{abstract}

\section{Introduction}
Session~B follows Session~A planning (\texttt{SCOPE.md}, \texttt{ARCHITECTURE.md}) by replacing scaffold stubs with working modules under \texttt{src/} and a four-tab Streamlit UI in \texttt{app/main.py}.

\section{Objectives}
\begin{itemize}[noitemsep]
  \item Implement data loading, algorithms, and UI (Exercise~1).
  \item Document AI interactions in \texttt{prompt\_log.md} (Exercise~2).
  \item Self-review tests and physical checks (Exercise~3).
\end{itemize}

\section{Exercise 1: Core features}

\subsection{Step 1: Data loading}
\texttt{src/weather/load.py} loads \texttt{data/sample\_rainfall.csv} with columns \texttt{datetime}, \texttt{rainfall\_mm}. Validation raises if columns are missing.

\subsection{Step 2: Core algorithms}
\begin{table}[htbp]
  \centering
  \caption{Implemented modules (reused from Weeks 3--6).}
  \label{tab:modules}
  \small
  \begin{tabular}{@{}ll@{}}
    \toprule
    Module & Function \\
    \midrule
    \texttt{src/weather/alerts.py} & GREEN / AMBER / RED vs.\ mm/h thresholds \\
    \texttt{src/runoff/scs\_cn.py} & SCS-CN runoff; enforce $Q \le P$ \\
    \texttt{src/reservoir/optimizer.py} & \texttt{scipy.optimize.minimize\_scalar} horizon \\
    \texttt{src/flood/inundation.py} & Wet mask, depth, \texttt{data/dem.npy} \\
    \bottomrule
  \end{tabular}
\end{table}

\textbf{Unit tests:} \texttt{tests/test\_runoff.py}, \texttt{tests/test\_flood.py}. Fixed \texttt{ModuleNotFoundError} via \texttt{pytest.ini} (\texttt{pythonpath = .}) and \texttt{conftest.py}.

\textbf{Terminal:} \texttt{pytest -q} $\rightarrow$ \texttt{2 passed in 0.10s}.

\subsection{Step 3: Streamlit UI}
\texttt{streamlit run app/main.py} exposes four tabs: Weather \& Alerts, Runoff (SCS-CN), Reservoir, Flood map. Figure~\ref{fig:streamlit} shows the Weather tab: peak 22~mm/h triggers a \textbf{RED} alert with default red threshold 20~mm/h.

\begin{figure}[htbp]
  \centering
  \appfig{week7_session_b_Streamlit_page.png}
  \caption{Streamlit dashboard (Weather \& Alerts tab) at \texttt{http://localhost:8501}.}
  \label{fig:streamlit}
\end{figure}

\section{Exercise 2: Prompt log}
Session~B entries in \texttt{prompt\_log.md} record the pytest path fix, module wiring prompts, run commands, GitHub push, and screenshot filename. Full file in Appendix~\ref{app:promptlog}.

\section{Exercise 3: Code review}
\begin{itemize}[noitemsep]
  \item Runoff satisfies $Q \le P$ for tested storms.
  \item Reservoir schedule matches Week~6 Lab~3 revenue and PASS validation.
  \item Flood uses \texttt{elevation <= flood\_level}; DEM bundled in repo.
  \item No API keys or tokens committed.
\end{itemize}

\section{Exercise 4: Repository update}
\textbf{URL:} \url{\capstonerepo}

\textbf{Commit:} \texttt{Week 7 Session B: Streamlit dashboard and pytest fix} (\texttt{408f211})

Session~A planning commits remain on \texttt{main}; Session~B adds \texttt{app/main.py}, \texttt{src/} implementations, \texttt{data/dem.npy}, and test configuration.

\section{Discussion}
\begin{enumerate}[noitemsep]
  \item Reusing lab code in \texttt{src/} packages kept the one-hour lab feasible.
  \item \texttt{pythonpath} configuration is required for both \texttt{pytest} and Streamlit imports.
  \item One dashboard screenshot plus passing tests meets optional submission evidence.
\end{enumerate}

\section{Lessons learned}
Configure \texttt{pytest.ini} when using a \texttt{src/} layout. Restart Streamlit after replacing stub tabs. Push to GitHub after each lab session; document commits in \texttt{prompt\_log.md}.

\section{Conclusion}
Week~7 Session~B objectives were met: working capstone dashboard, unit tests passing, code on GitHub, and documented AI collaboration. Next: Week~8 testing/validation and presentation.

\appendix

\section{Streamlit entry: \texttt{app/main.py}}
\label{app:main}
\filebox{\texttt{app/main.py} (full file)}{main.py}

\section{Runoff module: \texttt{src/runoff/scs\_cn.py}}
\label{app:runoff}
\filebox{\texttt{runoff\_scs\_cn.py} (excerpt file)}{runoff_scs_cn.py}

\section{Submitted prompt log: \texttt{prompt\_log.md}}
\label{app:promptlog}
\filebox{\texttt{prompt\_log.md} (full file, Sessions A + B)}{prompt_log.md}

\end{document}
